% Part: proof-theory
% Chapter: normalization
% Section: permutations

\documentclass[../../../include/open-logic-section]{subfiles}

\begin{document}

\olfileid{pt}{nor}{per}

\olsection{ક્રમપરિવર્તન-રૂપાંતરો}

\Log{N1i}માં કટ
એવા ખંડો છે
જેમનો અંત
!!a{elimination} નિયમના
મુખ્ય આધારવિધાનમાં
થાય છે. સામાન્યીકરણની
સાબિતીમાં એક કિસ્સો
મહત્તમ કટની લંબાઈ
અને તેની સાથે
આખી !!{proof}ની
કટ-લંબાઈ ઘટાડવાનો છે.
$>1$ લંબાઈનો કટ
ઘણી રીતે પૂરો
થઈ શકે છે: કટમાં
$!A$ની છેલ્લી
!!{formula} આવૃત્તિ
(\Elim{\lor} કે
\Elim{\lexists})નું
નિષ્કર્ષ છે કે નહીં,
અને તે કયા
!!{elimination} અનુમાનનું
મુખ્ય આધારવિધાન
છે, તેના પર
તે આધાર રાખે છે.

દાખલા તરીકે, સંબંધિત
અનુમાનો~\Elim{\lor}
અને~\Elim{\land}
હોય, $!A \ident !B \land !C$
હોય, અને
ઉપ-!!{proof}~$\delta_1$નો
અંત આ રીતે
થતો હોય:
\sourcecorrection{OLMD-253}{મૂળ ગદ્યમાં $!A$ને $!B \lor !C$ કહે છે, પરંતુ નીચેના વૃક્ષમાં કટનું છેલ્લું મુખ્ય સૂત્ર $!B \land !C$ છે; દાખલાના અનુરૂપ અહીં સંયોજન રાખ્યું છે.}
\[
\AxiomC{}
\RightLabel{$\delta_2$}
\DeduceC{$!D \lor !E$}
\AxiomC{$\Discharge{!D}{x}$}
\RightLabel{$\delta_3$}
\DeduceC{$!B \land !C$}
\AxiomC{$\Discharge{!E}{y}$}
\RightLabel{$\delta_4$}
\DeduceC{$!B \land !C$}
\DischargeRule{\Elim{\lor}}{x\,y}
\TrinaryInfC{$!B \land !C$}
\RightLabel{\Elim{\land}[1]}
\UnaryInfC{$!B$}
\DisplayProof
\]
કટમાં છેલ્લી
!!{formula} આવૃત્તિ~$!A_n$
એ~\Elim{\lor}નું નિષ્કર્ષ,
નીચેનું $!B \land !C$,
છે. છેલ્લાથી
આગળની
!!{formula} આવૃત્તિ~$!A_{n-1}$
એ~\Elim{\lor}ના
ગૌણ આધારવિધાનોમાંથી
એક છે.

આ કટની લંબાઈ
ઘટાડવા~\Elim{\lor}
અને~\Elim{\land}
અનુમાનોનો ક્રમ
બદલી (``ક્રમપરિવર્તન''
કરી) શકાય છે.
એટલે~$\delta$માંની
ઉપસાબિતી~$\delta_1$ના
સ્થાને આ મૂકીએ:
\[
\AxiomC{}
\RightLabel{$\delta_2$}
\DeduceC{$!D \lor !E$}
\AxiomC{$\Discharge{!D}{x}$}
\RightLabel{$\delta_3$}
\DeduceC{$!B \land !C$}
\RightLabel{\Elim{\land}[1]}
\UnaryInfC{$!B$}
\AxiomC{$\Discharge{!E}{y}$}
\RightLabel{$\delta_4$}
\DeduceC{$!B \land !C$}
\RightLabel{\Elim{\land}[1]}
\UnaryInfC{$!B$}
\DischargeRule{\Elim{\lor}}{x\,y}
\TrinaryInfC{$!B$}
\DisplayProof
\]
અગાઉ~\Elim{\lor}ની
નીચેના~\Elim{\land}ના
આધારવિધાનમાં પૂરા
થતા કટો હવે
\Elim{\lor}ના ગૌણ
આધારવિધાનોની ઉપરના
\Elim{\land} અનુમાનોમાંથી
એકના આધારવિધાનમાં
પૂરા થાય છે.
એટલે દરેકની
લંબાઈ~$1$ ઘટે છે.

વળી, સંપૂર્ણપણે
$\delta_2$, $\delta_3$
અથવા~$\delta_4$ની
અંદરનો કટ, કે
$\delta$માં~$\delta_1$ની
સંપૂર્ણપણે બહારનો
કટ, યથાવત્ રહે
છે: તે જ
!!{formula}s અને
તે જ અનુમાનોમાંથી
પસાર થાય છે.
તેથી ખાસ કરીને,
$\delta$માંના અનુરૂપ
કટ જેટલી જ તેની
કક્ષા અને લંબાઈ
રહે છે. આમ નવી
!!{proof}ની કટ-કક્ષા
વધતી નથી અને
નવી !!{proof}ની
કટ-લંબાઈ
જૂની !!{proof}ની
કટ-લંબાઈ કરતાં
ઓછી છે.

\begin{prob}
ક્રમપરિવર્તન-રૂપાંતર પછી
ઓછામાં ઓછા એક
કટની લંબાઈ~$1$
ઘટે છે.
\Elim{\land}ને~\Elim{\lor}
ની આરપાર ખસેડીએ
ત્યારે હકીકતમાં
બે કટ-ખંડોની
લંબાઈ ઘટે છે;
એથી આ કિસ્સામાં
!!{proof}ની કટ-લંબાઈ
ઓછામાં ઓછી~$2$
ઘટે છે. એક જ
આવા ક્રમપરિવર્તનથી
!!a{proof}ની કટ-લંબાઈ
$2$થી વધુ ઘટી
શકે? કટની
\emph{કક્ષા} ઘટી શકે?
\end{prob}

!!a{elimination} નિયમને
\Elim{\lor} કે
\Elim{\lexists} નિયમની
ઉપર ખસેડવાની
પ્રક્રિયાને
\emph{ક્રમપરિવર્તન-રૂપાંતર}
કહે છે. કેટલાક
નિયમ-યુગ્મો માટેનાં
નમૂના
\olref{tab:permutation-conversions}માં
આપ્યાં છે.
(બાકીનાં કસરત
તરીકે મૂક્યાં છે.)

% \begin{table}
% \begin{multline*}
% \shoveleft{\AxiomC{$!D \lor !E$}
% \AxiomC{$\Discharge{!D}{x}$}
% \shortDeduce
% \DeduceC{$!B \land !C$}
% \AxiomC{$\Discharge{!E}{y}$}
% \shortDeduce
% \DeduceC{$!B \land !C$}
% \DischargeRule{\Elim{\lor}}{x\,y}
% \TrinaryInfC{$!B \land !C$}
% \RightLabel{\Elim{\land}[1]}
% \UnaryInfC{$!B$}
% \DisplayProof}\\
% \shoveright{\longrightarrow\ 
% \AxiomC{$!D \lor !E$}
% \AxiomC{$\Discharge{!D}{x}$}
% \shortDeduce
% \DeduceC{$!B \land !C$}
% \RightLabel{\Elim{\land}[1]}
% \UnaryInfC{$!B$}
% \AxiomC{$\Discharge{!E}{y}$}
% \shortDeduce
% \DeduceC{$!B \land !C$}
% \RightLabel{\Elim{\land}[1]}
% \UnaryInfC{$!B$}
% \DischargeRule{\Elim{\lor}}{x\,y}
% \TrinaryInfC{$!B$}
% \DisplayProof}
% \\
% \shoveleft{\AxiomC{$!D \lor !E$}
% \AxiomC{$\Discharge{!D}{x}$}
% \shortDeduce
% \DeduceC{$!B \lif !C$}
% \AxiomC{$\Discharge{!E}{y}$}
% \shortDeduce
% \DeduceC{$!B \lif !C$}
% \DischargeRule{\Elim{\lor}}{x\,y}
% \TrinaryInfC{$!B \lif !C$}
% \AxiomC{$!B$}
% \RightLabel{\Elim{\lif}}
% \BinaryInfC{$!C$}
% \DisplayProof}\\
% \shoveright{\longrightarrow\ 
% \AxiomC{$!D \lor !E$}
% \AxiomC{$\Discharge{!D}{x}$}
% \shortDeduce
% \DeduceC{$!B \lif !C$}
% \AxiomC{$!B$}
% \RightLabel{\Elim{\lif}}
% \BinaryInfC{$!C$}
% \AxiomC{$\Discharge{!E}{y}$}
% \shortDeduce
% \DeduceC{$!B \lif !C$}
% \AxiomC{$!B$}
% \RightLabel{\Elim{\lif}}
% \BinaryInfC{$!C$}
% \DischargeRule{\Elim{\lor}}{x\,y}
% \TrinaryInfC{$!C$}
% \DisplayProof}
% \\
% \shoveleft{\AxiomC{$!D \lor !E$}
% \AxiomC{$\Discharge{!D}{x}$}
% \shortDeduce
% \DeduceC{$!B \lif !C$}
% \AxiomC{$\Discharge{!E}{y}$}
% \shortDeduce
% \DeduceC{$!B \lif !C$}
% \DischargeRule{\Elim{\lor}}{x\,y}
% \TrinaryInfC{$!B \lif !C$}
% \AxiomC{$!B$}
% \RightLabel{\Elim{\lif}}
% \BinaryInfC{$!C$}
% \DisplayProof}\\
% \shoveright{\longrightarrow\ 
% \AxiomC{$!D \lor !E$}
% \AxiomC{$\Discharge{!D}{x}$}
% \shortDeduce
% \DeduceC{$!B \lif !C$}
% \AxiomC{$!B$}
% \RightLabel{\Elim{\lif}}
% \BinaryInfC{$!C$}
% \AxiomC{$\Discharge{!E}{y}$}
% \shortDeduce
% \DeduceC{$!B \lif !C$}
% \AxiomC{$!B$}
% \RightLabel{\Elim{\lif}}
% \BinaryInfC{$!C$}
% \DischargeRule{\Elim{\lor}}{x\,y}
% \TrinaryInfC{$!C$}
% \DisplayProof}
% \\
% \shoveleft{\AxiomC{$!D \lor !E$}
% \AxiomC{$\Discharge{!D}{x}$}
% \shortDeduce
% \DeduceC{$\lforall[x][!B(x)]$}
% \AxiomC{$\Discharge{!E}{y}$}
% \shortDeduce
% \DeduceC{$\lforall[x][!B(x)]$}
% \DischargeRule{\Elim{\lor}}{x\,y}
% \TrinaryInfC{$\lforall[x][!B(x)]$}
% \RightLabel{\Elim{\lforall}}
% \UnaryInfC{$!B(t)$}
% \DisplayProof}\\
% \shoveright{\longrightarrow\ 
% \AxiomC{$!D \lor !E$}
% \AxiomC{$\Discharge{!D}{x}$}
% \shortDeduce
% \DeduceC{$\lforall[x][!B(x)]$}
% \RightLabel{\Elim{\lforall}}
% \UnaryInfC{$!B(t)$}
% \AxiomC{$\Discharge{!E}{y}$}
% \shortDeduce
% \DeduceC{$\lforall[x][!B(x)]$}
% \RightLabel{\Elim{\lforall}}
% \UnaryInfC{$!B(t)$}
% \DischargeRule{\Elim{\lor}}{x\,y}
% \TrinaryInfC{$!B(t)$}
% \DisplayProof}
% \end{multline*}
% \caption{Some permutation conversions for \Log{N1i}.}
% \ollabel{tab:permutation-conversions}
% \end{table}

{\small
\begin{longtable}{@{}l@{}}
\AxiomC{$!D \lor !E$}
\AxiomC{$\Discharge{!D}{x}$}
\shortDeduce
\DeduceC{$!B \land !C$}
\AxiomC{$\Discharge{!E}{y}$}
\shortDeduce
\DeduceC{$!B \land !C$}
\DischargeRule{\Elim{\lor}}{x\,y}
\TrinaryInfC{$!B \land !C$}
\RightLabel{\Elim{\land}[1]}
\UnaryInfC{$!B$}
\DisplayProof
$\longrightarrow$\\[6ex]
\quad\AxiomC{$!D \lor !E$}
\AxiomC{$\Discharge{!D}{x}$}
\shortDeduce
\DeduceC{$!B \land !C$}
\RightLabel{\Elim{\land}[1]}
\UnaryInfC{$!B$}
\AxiomC{$\Discharge{!E}{y}$}
\shortDeduce
\DeduceC{$!B \land !C$}
\RightLabel{\Elim{\land}[1]}
\UnaryInfC{$!B$}
\DischargeRule{\Elim{\lor}}{x\,y}
\TrinaryInfC{$!B$}
\DisplayProof
\\[10ex]
\AxiomC{$!D \lor !E$}
\AxiomC{$\Discharge{!D}{x}$}
\shortDeduce
\DeduceC{$!B \lif !C$}
\AxiomC{$\Discharge{!E}{y}$}
\shortDeduce
\DeduceC{$!B \lif !C$}
\DischargeRule{\Elim{\lor}}{x\,y}
\TrinaryInfC{$!B \lif !C$}
\AxiomC{$!B$}
\RightLabel{\Elim{\lif}}
\BinaryInfC{$!C$}
\DisplayProof
$\longrightarrow$\\[6ex]
\AxiomC{$!D \lor !E$}
\AxiomC{$\Discharge{!D}{x}$}
\shortDeduce
\DeduceC{$!B \lif !C$}
\AxiomC{$!B$}
\RightLabel{\Elim{\lif}}
\BinaryInfC{$!C$}
\AxiomC{$\Discharge{!E}{y}$}
\shortDeduce
\DeduceC{$!B \lif !C$}
\AxiomC{$!B$}
\RightLabel{\Elim{\lif}}
\BinaryInfC{$!C$}
\DischargeRule{\Elim{\lor}}{x\,y}
\TrinaryInfC{$!C$}
\DisplayProof
\\[10ex]
\AxiomC{$!D \lor !E$}
\AxiomC{$\Discharge{!D}{x}$}
\shortDeduce
\DeduceC{$!B \lif !C$}
\AxiomC{$\Discharge{!E}{y}$}
\shortDeduce
\DeduceC{$!B \lif !C$}
\DischargeRule{\Elim{\lor}}{x\,y}
\TrinaryInfC{$!B \lif !C$}
\AxiomC{$!B$}
\RightLabel{\Elim{\lif}}
\BinaryInfC{$!C$}
\DisplayProof
$\longrightarrow$\\[6ex]
\AxiomC{$!D \lor !E$}
\AxiomC{$\Discharge{!D}{x}$}
\shortDeduce
\DeduceC{$!B \lif !C$}
\AxiomC{$!B$}
\RightLabel{\Elim{\lif}}
\BinaryInfC{$!C$}
\AxiomC{$\Discharge{!E}{y}$}
\shortDeduce
\DeduceC{$!B \lif !C$}
\AxiomC{$!B$}
\RightLabel{\Elim{\lif}}
\BinaryInfC{$!C$}
\DischargeRule{\Elim{\lor}}{x\,y}
\TrinaryInfC{$!C$}
\DisplayProof
\\[10ex]
\AxiomC{$!D \lor !E$}
\AxiomC{$\Discharge{!D}{x}$}
\shortDeduce
\DeduceC{$\lforall[x][!B(x)]$}
\AxiomC{$\Discharge{!E}{y}$}
\shortDeduce
\DeduceC{$\lforall[x][!B(x)]$}
\DischargeRule{\Elim{\lor}}{x\,y}
\TrinaryInfC{$\lforall[x][!B(x)]$}
\RightLabel{\Elim{\lforall}}
\UnaryInfC{$!B(t)$}
\DisplayProof
\quad$\longrightarrow$\\[8ex]
\quad\AxiomC{$!D \lor !E$}
\AxiomC{$\Discharge{!D}{x}$}
\shortDeduce
\DeduceC{$\lforall[x][!B(x)]$}
\RightLabel{\Elim{\lforall}}
\UnaryInfC{$!B(t)$}
\AxiomC{$\Discharge{!E}{y}$}
\shortDeduce
\DeduceC{$\lforall[x][!B(x)]$}
\RightLabel{\Elim{\lforall}}
\UnaryInfC{$!B(t)$}
\DischargeRule{\Elim{\lor}}{x\,y}
\TrinaryInfC{$!B(t)$}
\DisplayProof
\\
\caption{\Log{N1i} માટેનાં કેટલાંક ક્રમપરિવર્તન-રૂપાંતરો.}
\ollabel{tab:permutation-conversions}
\end{longtable}
}
\sourcecorrection{OLMD-254}{મૂળ કોષ્ટકમાં અનુસરણ-લોપના બે મધ્યવર્તી રૂપાંતર-નમૂના એકસરખા પુનરાવર્તિત થાય છે. મૂળના વૃક્ષો યથાવત્ રાખ્યાં છે; કોઈ ગાયબ નમૂનો અનુમાનથી ઉમેર્યો નથી.}

\begin{prob}
\Elim{\lor} પછી
\Elim{\lor} અથવા
\Elim{\lnot} આવે
ત્યારેનાં
ક્રમપરિવર્તન-રૂપાંતરો
આપો.
\end{prob}

\begin{prob}
\Elim{\lexists} પછી
\Elim{\exists} આવે
ત્યારેનાં
ક્રમપરિવર્તન-રૂપાંતરો
આપો. છેલ્લા
કિસ્સામાં કોઈ
સ્વનિર્ધારક ચલની
શરત કેમ ભંગાતી
નથી તે સમજાવો.
\end{prob}

ક્રમપરિવર્તન-રૂપાંતર
લગાવતાં કોઈ
સ્વનિર્ધારક ચલની
શરત ભંગાતી નથી,
તેની ખાતરી કરીએ.
આ જુઓ:
\[
\AxiomC{}
\RightLabel{$\delta_2$}
\DeduceC{$!D \lor !E$}
\AxiomC{$\Discharge{!D}{x}$}
\RightLabel{$\delta_3$}
\DeduceC{$\lexists[x][!B(x)]$}
\AxiomC{$\Discharge{!E}{y}$}
\RightLabel{$\delta_4$}
\DeduceC{$\lexists[x][!B(x)]$}
\DischargeRule{\Elim{\lor}}{x\,y}
\TrinaryInfC{$\lexists[x][!B(x)]$}
\AxiomC{$\Discharge{!B(c)}{z}$}
\RightLabel{$\delta_5$}
\DeduceC{$!C$}
\DischargeRule{\Elim{\lexists}}{z}
\BinaryInfC{$!C$}
\DisplayProof
\]
તેના સ્થાને
આ મૂકીએ:
\[
\AxiomC{}
\RightLabel{$\delta_2$}
\DeduceC{$!D \lor !E$}
\AxiomC{$\Discharge{!D}{x}$}
\RightLabel{$\delta_3$}
\DeduceC{$\lexists[x][!B(x)]$}
\AxiomC{$\Discharge{!B(c)}{z}$}
\RightLabel{$\delta_5$}
\DeduceC{$!C$}
\DischargeRule{\Elim{\lexists}}{z}
\BinaryInfC{$!C$}
\AxiomC{$\Discharge{!E}{y}$}
\RightLabel{$\delta_4$}
\DeduceC{$\lexists[x][!B(x)]$}
\AxiomC{$\Discharge{!B(c)}{z}$}
\RightLabel{$\delta_5$}
\DeduceC{$!C$}
\DischargeRule{\Elim{\lexists}}{z}
\BinaryInfC{$!C$}
\DischargeRule{\Elim{\lor}}{x\,y}
\TrinaryInfC{$!C$}
\DisplayProof
\]
તો કદાચ
સ્વનિર્ધારક ચલ~$c$
દાખલા તરીકે~$!D$માં
આવતું હોય એવી
ચિંતા થાય. મૂળ
વૃક્ષમાં ધારણા~$!D$
મૂળના~$\Elim{\lexists}$ની
ઉપર મુક્ત થાય છે.
તેથી મૂળ !!{proof}માં
તે~\Elim{\exists}
અનુમાનના મુખ્ય
આધારવિધાનની ઉપર
ખુલ્લી રહેલી
ધારણામાં આવતું
નથી; ત્યાં
સ્વનિર્ધારક ચલની
શરત સંતોષાય છે.
પણ નવી !!{proof}માં
$!D$~\Elim{\lexists}
અનુમાનો પછી જ
મુક્ત થાય છે,
એટલે શરત ભંગાતી
હોય એમ લાગે.
જોકે યાદ રાખો
કે આપણે
!!{proof}s નિયમિત
છે એમ માનીએ છીએ:
દરેક સ્વનિર્ધારક
ચલ માત્ર એક
અનુમાનનું સ્વનિર્ધારક
ચલ છે અને આખી
!!{proof}માં તે
માત્ર~\Elim{\lexists}ના
ગૌણ આધારવિધાનની
ઉપર જ આવે છે.
ખાસ કરીને,
$c$~$\delta_3$માં
કે~$\delta_4$માં
આવી શકતું નથી.

આ દાખલો બીજી
એક વાત પણ
બતાવે છે: નવી
!!{proof}માં~$\delta_5$ની
બે નકલો છે.
જો~$\delta_5$માં
મહત્તમ કટ હોય,
તો આપણે કટ-લંબાઈ
\emph{ઘટાડી નથી}!
તેથી~$\delta_5$માં
મહત્તમ કટ નથી
તેની ખાતરી હોય
ત્યારે જ
ક્રમપરિવર્તન-રૂપાંતર
લગાવવું જોઈએ.
આ માટે હંમેશાં
\emph{સૌથી ઉપરનો}
મહત્તમ કટ પસંદ
કરીશું: એટલે
એવો મહત્તમ કટ
જેની ઉપ-!!{proof}
(અહીં~$\delta_1$)માં
તે ઉપ-!!{proof}ના
છેલ્લા અનુમાનની
ઉપર પૂરો થતો
બીજો કોઈ મહત્તમ
કટ ન હોય.
એથી~$\delta_5$માં
(અને ખરેખર
$\delta_2$, $\delta_3$
કે~$\delta_4$માં પણ)
બીજા મહત્તમ
કટ હોવાની શક્યતા
દૂર થાય છે.

\end{document}
